\subsubsection{Examples of Operators}\label{subsec:operator-examples}

\Needspace{4\baselineskip}
\begin{Example}
\label{ex:dual-functional-bounded-operator}
For any Banach space $X$, any element of the dual space $X^*$ is by definition in $\mathcal{L}(X,\mathbb{R})$.
\end{Example}

\Needspace{4\baselineskip}
\begin{Example}
\label{ex:identity-operator-and-inclusion}
Of course, the \textbf{identity operator} $I$ from $X$ to $X$ defined by $Ix=x$ is a trivial example of a bounded operator. More interestingly, if $X\subset Y$, we can define an identity mapping from $X$ to $Y$ to be the map that takes each $x\in X$ to the same $x$ (though this time thought of as an element of $Y$).
\end{Example}

The following lemma is an immediate consequence of Definition~\ref{def:compact-imbedding}.

\begin{Lemma}
\label{lem:identity-bounded-continuous-embedding}
Let $X$ and $Y$ be Banach spaces. Then $X$ is continuously imbedded in $Y$ if and only if the identity mapping from $X$ to $Y$ is well defined (i.e., $X\subset Y$) and bounded.
\end{Lemma}

We use this in the following example.

\Needspace{4\baselineskip}
\begin{Example}
\label{ex:sobolev-identity-inclusion-bounded}
Corollary~\ref{cor:sobolev-extension-continuity} implies that if $k>m/2$ and $\Omega\subset\mathbb{R}^m$ has the $k$-extension property, then the identity mapping from $H^k(\Omega)$ to $C_b(\overline\Omega)$ is bounded.
\end{Example}

\Needspace{4\baselineskip}
\begin{Example}
\label{ex:riesz-map-conjugate-linear-operator}
The Riesz representation theorem states that every Hilbert space $H$ is isometric to its dual space $H^*$. In Remark~\ref{rem:riesz-mappings} we constructed the map $A_H:H^*\to H$ that, for every linear functional $l\in H^*$, gives the unique element $x=A_H(l)\in H$ such that\[l(y)=(x,y)\quad\text{for all }y\in H.\]Note that the operator $A_H$ is actually \emph{conjugate linear}; i.e., $A_H(\alpha l)=\bar\alpha A_H(l)$. The operator is bounded; the Riesz representation theorem assures us that its norm is $1$.
\end{Example}

\Needspace{4\baselineskip}
\begin{Example}
\label{ex:right-left-shifts-domain-dependent-inverses}
For any $x=\{x_1,x_2,\ldots,x_n,\ldots\}\in\ell^2$ we define the \textbf{right shift} operator $S_r:\ell^2\to\ell^2$ by\[S_r(x)=\{0,x_1,x_2,\ldots,x_n,\ldots\}.\]The range is the set of all elements in $\ell^2$ with first component $0$; i.e., $\{1,0,0,\ldots\}^{\perp}$. The nullspace is the singleton $\{0\}$. Since $\|S_r(x)\|=\|x\|$, we have $\|S_r\|=1$.

We define the \textbf{left shift} operator $S_l:\ell^2\to\ell^2$ by\[S_l(x)=\{x_2,\ldots,x_n,\ldots\}.\]The range of $S_l$ is $\ell^2$ and its nullspace is spanned by the element $\{1,0,0,\ldots\}$. Since $\|S_l(x)\|\leq\|x\|$ and since $\|S_l(x)\|=\|x\|$ for any $x\in\{1,0,0,\ldots\}^{\perp}$ we have $\|S_l\|=1$.

Clearly, $S_r$ is invertible and $S_r^{-1}=S_l$. If we take $\mathcal{D}(S_l)=\ell^2$, then $S_l$ is not invertible. (Its nullspace is nontrivial.) However, if we take $\mathcal{D}(S_l):=\{1,0,0,\ldots\}^{\perp}$, then $S_l$ is invertible and $S_l^{-1}=S_r$.
\end{Example}

\Needspace{4\baselineskip}
\begin{Example}
\label{ex:orthogonal-projection-operator}
Let $H$ be a Hilbert space and let $M\subset H$ be a closed subspace. Then we define the \textbf{orthogonal projection} operator $P_M:H\to M$ so that for any $x\in H$, $P_Mx$ is the unique element of $M$ such that $(x-P_Mx)\in M^{\perp}$.
\end{Example}
